\chapter*{\centering ABSTRAK}
\addcontentsline{toc}{chapter}{ABSTRAK}

\begin{center}
\setstretch{1.2}
\begin{tabular}{ll}
Nama          & : Fahry Rizky Samsudin \\ 
Program Studi & : Fisika \\ 
Judul         & : \begin{tabular}[t]{@{}l@{}}Rancang Bangun Sistem Visi Komputer Tertanam Berbasis\\Raspberry Pi 5 untuk Kendali Robot 4-DOF: Teleoperasi\\Dobot Magician dan Pengenalan SIBI pada meArm\end{tabular} \\ 
\end{tabular}
\end{center}

\vspace{0.8cm}
\begin{spacing}{1.25}
\noindent Penelitian ini berfokus pada perancangan sistem tertanam berbasis Raspberry Pi 5 untuk kendali dua platform robot manipulator 4-DOF, yaitu Dobot Magician dan meArm. Sistem dirancang dalam dua mode kerja: mode teleoperasi untuk mengendalikan gerak Dobot Magician berdasarkan pelacakan tangan, serta mode pengenalan Sistem Isyarat Bahasa Indonesia (SIBI) untuk memetakan isyarat tangan menjadi respons gerak pada meArm. Metode yang digunakan meliputi studi literatur, integrasi perangkat keras, pengembangan piranti lunak berbasis Python 3.9, ekstraksi landmark tangan menggunakan MediaPipe, pemetaan koordinat citra ke ruang kerja robot, serta pengujian akurasi klasifikasi, presisi posisi, dan waktu respons. Laporan ini diharapkan menjadi dasar pengembangan sistem kendali robot yang ringan, portabel, dan mampu berjalan pada perangkat komputasi tepi untuk kebutuhan edukasi robotika, interaksi manusia--robot, dan pengenalan bahasa isyarat.

\vspace{0.5cm}
\noindent\textbf{Kata kunci:} Raspberry Pi 5, sistem tertanam, MediaPipe, SIBI, robot manipulator 4-DOF.
\end{spacing}
\newpage
